\chapterhead{5}{The Long Loop}{Following water from sea to sky and back}\label{ch5}

Evaporation, condensation and a place to flow back make a loop. On a huge scale, water evaporates from the sea and from wet land, rises as vapour, cools and condenses into tiny drops that make clouds, falls as rain or snow, and runs back toward the sea. That is the \textbf{water cycle}.

\begin{center}
\begin{tikzpicture}[font=\sffamily\scriptsize]
  \node[bigcir,minimum size=1.9cm] (sea) at (0,0) {sea};
  \node[bigcir,minimum size=1.9cm] (air) at (3.4,2.0) {air};
  \node[bigcir,minimum size=1.9cm] (cloud) at (7.0,2.0) {cloud};
  \node[bigcir,minimum size=1.9cm] (rain) at (9.6,0) {rain};
  \node[bigcir,minimum size=1.9cm] (river) at (5.0,-1.8) {river};
  \draw[moveA] (sea) -- node[above left,note] {evaporation} (air);
  \draw[moveA] (air) -- node[above,note] {condensation} (cloud);
  \draw[moveA] (cloud) -- node[above right,note] {falls} (rain);
  \draw[moveA] (rain) -- node[below right,note] {flows} (river);
  \draw[moveA] (river) -- node[below left,note] {flows} (sea);
\end{tikzpicture}
\end{center}

\begin{tryit}{A solar still}
\textbf{You need:} a large bowl, a small cup, cling film, tape, a small stone, water, food colouring, a spoon, a sunny window sill.
\begin{enumerate}[leftmargin=1.6em]
\item Put 20 spoonfuls of coloured water in the bowl. Stand the small cup in the middle. Do not let water get in the cup.
\item Cover the bowl with cling film, taped tight. Put the small stone on the film right over the cup so the film dips toward it.
\item Put the bowl in the sun for a day. Look for drops on the film. Where do they run to?
\item At the end, count spoonfuls in the cup, in the bowl, and note any drops left on the film. Is the cup's water coloured?
\end{enumerate}
\end{tryit}
\adult{Cold tap water only. Do not drink any water from this activity. Put the bowl where it will not be knocked.}

\begin{worked}{Practice numbers, not from a real test}
Start: 20 spoonfuls of coloured water.\\
End: 16 in the bowl, 3 in the cup, and about 1 spoonful worth still as drops on the film.\\
$16+3+1=20$. The water in the cup was clear. The colour stayed behind in the bowl.
\end{worked}

\begin{center}
\begin{tikzpicture}
  \node[cir] (a) at (0,0) {20};
  \node[cir] (b) at (3.6,1.2) {16}; \node[cir] (c) at (3.6,0) {3}; \node[cir] (d) at (3.6,-1.2) {1};
  \draw[moveA] (a) -- (b); \draw[moveA] (a) -- (c); \draw[moveA] (a) -- (d);
  \node[note,anchor=west] at (4.3,1.2) {left in the bowl};
  \node[note,anchor=west] at (4.3,0) {collected in the cup};
  \node[note,anchor=west] at (4.3,-1.2) {still on the film};
\end{tikzpicture}
\end{center}

\begin{multix}[hmm]\multixbot[hmm]\tcblower
The clear water in my cup is new. Clouds must make new water out of nothing. How else would rain appear? I like my idea. Do you?
\end{multix}

\trust{We counted every spoonful at the start and the end, including the drops on the film, and we checked the cup's water for colour.}
{The amounts added up to the start total, and the clear water in the cup shows the water travelled as vapour and left the colour behind.}
{Whether the real sea and sky work exactly the same way. A bowl in a window is a small stand-in for something enormously bigger.}

\begin{keepbox}
We turned the bowl into a small water cycle. \textbf{We kept} \wline[3.4cm] \textbf{and we set aside} \wline[3.4cm].
\end{keepbox}

\clearpage

\subhead{Rain in a bag}
\begin{tryit}{A cloud in a bag}
Put a spoonful of water in a zip bag, seal it, tape it to a sunny window and watch for a morning. Where do drops form? \wline[4cm]\\
What happens when drops get big? \wline[4cm]
\end{tryit}

\begin{worked}{Practice result}
After two hours small drops lined the top of the bag, where it was cooler. By the afternoon bigger drops ran down to the bottom. The water went round a tiny loop, much like the water cycle.
\end{worked}

\subhead{What a diagram of the cycle leaves out}
Every diagram of the water cycle is a model, a simple stand-in. Here is what the arrows above keep and what they set aside.

\begin{center}
\sffamily\small\renewcommand{\arraystretch}{1.5}
\begin{tabular}{@{}p{3.4cm}p{3.7cm}p{3.7cm}@{}}
\toprule
\textbf{Part of the diagram} & \textbf{What it keeps} & \textbf{What it sets aside}\\
\midrule
Boxes for the stages & the order of the stages & how many different places each is\\
Arrows between boxes & the direction water moves & how long each step takes\\
One loop & the same water going round & plants, ice, underground water\\
Labelled steps & the names of the changes & the amounts of water in each\\
\bottomrule
\end{tabular}
\end{center}

Water stays in the air for roughly a week or so on average. It can stay in the deep ocean or in ice for far longer. The diagram's arrows are all the same size, but the real steps are not.

\trust{We checked the model against the jar, the bowl and the bag, each of which keeps the total.}
{Each small test had the same order: liquid, vapour, drops, liquid again.}
{That a model with five boxes can predict rainfall, or how long water stays anywhere. It is a map of the order, not of the amounts or the times.}

\clearpage

\begin{puzzle}{Do the ledger}
A solar still started with 20 spoonfuls. At the end the bowl held 15 and the cup held 4. How many spoonfuls are unaccounted for? \wline[1.6cm]\\
Where might they be? \wline[4cm]
\end{puzzle}

\begin{puzzle}{Put them in order}
Number the steps from 1 to 5.\\
\wline[1cm] Rain falls. \quad \wline[1cm] Water evaporates from the sea. \quad \wline[1cm] Rivers flow to the sea.\\
\wline[1cm] Vapour condenses into cloud drops. \quad \wline[1cm] Vapour rises into the air.
\end{puzzle}

\begin{puzzle}{Which step?}
Write E for evaporation, C for condensation or R for rain.\\
(a) A puddle shrinks \wline[1cm] \quad (b) Dew on grass \wline[1cm] \quad (c) A cloud forms \wline[1cm]\\
(d) A wet street dries \wline[1cm] \quad (e) A steamy mirror \wline[1cm] \quad (f) A storm \wline[1cm]
\end{puzzle}

\begin{puzzle}{Clear water}
The water in the cup of the solar still was clear, but the water in the bowl was coloured. What does this tell you about what travelled? \wline[6cm]
\end{puzzle}

\begin{puzzle}{Follow a drop}
Start at the sea. Which of these routes is possible? Circle it.\\
\textbf{sea $\to$ air $\to$ cloud $\to$ rain $\to$ river $\to$ sea} \qquad \textbf{sea $\to$ cloud $\to$ river $\to$ air}
\end{puzzle}

\clearpage

\begin{puzzle}{Rain or snow?}
In cold weather the cloud drops may freeze. What might fall instead of rain? \wline[3cm]
\end{puzzle}

\begin{puzzle}{Missing from the diagram}
Name two things the five-box diagram of the water cycle leaves out. \wline[6cm]
\end{puzzle}

\begin{puzzle}{Same size arrows}
The diagram draws every arrow the same size. Why might that be misleading? \wline[6cm]
\end{puzzle}

\begin{puzzle}{Closed loop}
In the bag and the jar the water went round and round. Where in the real world is the water cycle \emph{not} closed in a small way? \wline[6cm]
\end{puzzle}

\begin{puzzle}{Your loop}
Draw your own loop for a puddle in the sun. Use circles and arrows. \wline[5cm]
\begin{center}\begin{tikzpicture}\draw[ink!70,line width=1pt,rounded corners=3pt] (0,0) rectangle (11,3);\end{tikzpicture}\end{center}
\end{puzzle}
\clearpage

\thispagestyle{fancy}
\subhead{Draw your own water cycle}
Draw a loop for a place you know, such as a lake, a garden or a rooftop. Use circles for the stages and arrows for the moves. Then fill in the three sentences.\par\medskip
\begin{center}\begin{tikzpicture}\draw[ink!70,line width=1pt,rounded corners=3pt] (0,0) rectangle (11.4,8.5);\end{tikzpicture}\end{center}
My diagram keeps: \wline[8cm]\\[10pt]
It sets aside: \wline[8cm]\\[10pt]
A question it cannot answer: \wline[6.5cm]
\clearpage
